% =============================================================================
%  Seccion: Conclusiones
%  Contenido de EJEMPLO.
% =============================================================================

\section{Conclusiones}\label{sec:conclusiones}

Que responden los resultados de \cref{sec:resultados} a la pregunta que se
planteo en \cref{sec:introduccion}, que limitaciones tiene el analisis y que
seguiria despues.

Un par de reglas para que el archivo compartido no se vuelva un campo minado de
conflictos de merge:

\begin{itemize}
  \item Escribe \textbf{una oracion por renglon} (o corta los renglones a ~80
    caracteres). Git compara linea por linea, y un parrafo en una sola linea
    gigante hace que cualquier cambio choque con el de tu companero.
  \item Agrega tus referencias nuevas hasta \textbf{el final} del archivo
    \texttt{references.bib}, sin reordenar las que ya estaban.
  \item No borres los \verb|\label{}|: alguien mas puede estar citandolos con
    \verb|\cref| desde otra seccion.
\end{itemize}

Las referencias aparecen automaticamente al final del documento: biblatex solo
imprime las entradas que de verdad se citaron. Si agregas algo a
\texttt{references.bib} y no lo citas, no sale en la lista.
